\documentclass[14pt,aspectratio=1610]{beamer}
%\documentclass[14pt,aspectratio=1610,handout]{beamer}
\usepackage{csu-theme}
\usepackage{arrays}
\usepackage{linked-lists}
\usetikzlibrary{positioning,backgrounds}

\title{\Huge{} Doubly-Linked Lists\\ \& Project 1}
\subtitle{CSCI 315: \textit{Data Structure Analysis}}
\author{Sean T. Hayes, Ph.D.}
\institute{Charleston Southern University}

%% Comment out date to use default (today)
% \date{February 11, 2025}
\subject{Data-Strucutre Analysis}
%\keywords{programming}

\setlength{\parskip}{0.5em}

\graphicspath{{../imgs/debugging}}

% Make the item in a description block have white text
\setbeamercolor{description item}{parent=normal text}
\begin{document}

\begin{frame}
  \titlepage{}
\end{frame}

\section{Motivation}

\begin{frame}{What if we want to traverse the list in reverse?}
  Three Options:\vspace{-0.25em}
\begin{enumerate}[<+->]\setlength\itemsep{0.5em}\setlength{\parskip}{0pt}
\item
  We can use recursion to iterate to the end and then perform the operation when we ``unwind.''
\item
  We can use a helper data structure (a stack).
\item
  We can implement a \emph{doubly-linked list} type, where every node has a next and previous pointer.
\end{enumerate}\onslide<4>{}%
\vspace{-0.25em}

\alt<-3>{
\begin{tikzpicture}[node of list/.append style = {minimum width=14.275mm}]
  \LinkPtr{mpHead}
  \LinkedNode[first]{"A"}
  \LinkedNode{"B"}
  \LinkedNode{\textellipsis}
  \LinkedNode{"Z"}
  \TailPtr{mpTail}
\end{tikzpicture}%
}{
\begin{tikzpicture}
  \LinkPtr{mpHead}
  \DoublyLinkedNode[first]{"A"}
  \DoublyLinkedNode{"B"}
  \DoublyLinkedNode{\textellipsis}
  \DoublyLinkedNode{"Z"}
  \TailPtr{mpTail}
\end{tikzpicture}%
}
\end{frame}

\section{Insertion}

\begin{frame}{Prepending a Doubly-Linked List}
\center
\begin{tikzpicture}


  \onslide<2->{
    \LinkPtr[-2]{pNewNode}
    \DoublyLinkedNode[first]{"new"}
    \node[block] (newNode) [fit=(ele)]{};
    \coordinate (newNodePrev) at (leftPtr);
    \coordinate (newNodeNext) at (ptr);
  }

  \LinkPtr{mpHead}
  \onslide<5->{
    \draw[link] (ptr) to[bend left=6] (newNode);
  }

  \alt<-4>{\DoublyLinkedNode[first]{"A"}}{\DoublyLinkedNode[first][disconnected]{"A"}}

  \onslide<3->{
    \draw[link] (newNodeNext) to [bend right=10] (ele);
    % \draw[link] (ptr) to[bend left=8] (ele);
  }
  \onslide<4->{
    \draw[link] (leftPtr) to [bend left=10] (newNode);
  }

  \DoublyLinkedNode{"B"}
  % \DoublyLinkedNode{"C"}
  \DoublyLinkedNode{\textellipsis}
  \DoublyLinkedNode{"Z"}
  \TailPtr{mpTail}
  \onslide<5>{}
\end{tikzpicture}
\end{frame}


\begin{frame}{Appending a Doubly-Linked List}
  \center%
  \begin{tikzpicture}

  
    \LinkPtr{mpHead}
    \DoublyLinkedNode[first]{"A"}
    \DoublyLinkedNode{"B"}
    \DoublyLinkedNode{\textellipsis}
    \DoublyLinkedNode{"Z"}
    \onslide<-4>{
      \TailPtr{mpTail}
    }
    \onslide<5->{
      \TailPtr[null]{mpTail}
    }
    \node[block] (oldTail) [fit=(ele)]{};%
    \coordinate (oldTailNext) at (ptr);%
    
      \onslide<2->{
        \begin{scope}[xshift=7.5cm]%
          \LinkPtr[2]{pNewNode}
        \end{scope}
        \DoublyLinkedNode[first]{"new"}
        \node[block] (newNode) [fit=(ele)]{};%
        \coordinate (newNodePrev) at (leftPtr);%
        \coordinate (newNodeNext) at (ptr);%
      }

      \onslide<3->{
        \draw[link] (leftPtr) to [bend right=8] (oldTail);%
      }

      \onslide<4->{
        \draw[link] (oldTailNext) to [bend right=8] (ele);%
      }

    
    \onslide<5>{
      \draw[link] (tPtr) to[bend right=60] (ele);%
    }
  \end{tikzpicture}
\end{frame}

\begin{frame}[fragile]{Insert in the Middle}
  
\begin{center}
  Example: Insert a new node after at\\ index 2 (after ``B'')

\begin{tikzpicture}
    \LinkPtr{mpHead}
    \DoublyLinkedNode[first]{"A"}
    \DoublyLinkedNode{"B"}
    \node[block] (BNode) [fit=(ele)]{};%
    \coordinate (beforeNodeNext) at (ptr);%
    \onslide<3->{
      \TailPtr{pBefore}
    }
    
    \alt<-5>{\DoublyLinkedNode{"C"}}{%
      \alt<-6>{\DoublyLinkedNode[disconnected]{"C"}}{%
        \DoublyLinkedNode[disconnected][disconnected]{"C"}%
      }%
    }
    \node[block] (afterNode) [fit=(ele)]{};%
    \coordinate (afterNodePrev) at (leftPtr);%
    \DoublyLinkedNode{"D"}
    \TailPtr{mpTail}

    \onslide<2->{
      \begin{scope}[xshift=3.75cm]%
        \LinkPtr[2]{pNewNode}
      \end{scope}
      \DoublyLinkedNode[first]{"new"}
      \node[block] (newNode) [fit=(ele)]{};%
      \coordinate (newNodePrev) at (leftPtr);%
      \coordinate (newNodeNext) at (ptr);%
    }

    \onslide<4->{
      \draw[link] (newNodePrev) to[bend right=8] (BNode);
    }
    \onslide<5->{
      \draw[link] (newNodeNext) to[bend left=8] (afterNode);
    }
    \onslide<6->{
      \draw[link] (afterNodePrev) to[bend left=8] (newNode);
    }
    \onslide<7->{
      \draw[link] (beforeNodeNext) to[bend right=8] (newNode);
    }
  \end{tikzpicture}
\end{center}
\end{frame}

\section{Faster Lookup}

\section{Remove}

\begin{frame}[fragile, c]{Remove an Element at an Index}
%\begin{center}
  \begin{tikzpicture}[node of list/.append style = {
    %minimum height=6mm,
    minimum width=3mm,
    node distance=6mm,
    %font=\small,
  },
  link pointer of list/.append style= {
    %minimum height=6mm,
    minimum width=3.5mm
  }]
    \LinkPtr{mpHead}
    \DoublyLinkedNode[first]{"A"}
    \DoublyLinkedNode{"B"}
    \alt<-3>{
      \DoublyLinkedNode{"C"}
    }{
      \onslide<-6>{
        \DoublyLinkedNode[middle][false]{"C"}
      }
    }
    \onslide<2->{
      \TailPtr{pToDel}
    }
    \alt<-5>{
      \DoublyLinkedNode[middle][connect]{"D"}
    }{
      \alt<-6>{
      \DoublyLinkedNode[false][connect]{"D"}
      }{
        \DoublyLinkedNode[false][false]{"D"}
      }
    }
    \onslide<4->{
      \draw[link] ({nextPtr"B"}) to[bend left] (ele);%
    }

    \onslide<6->{
      \draw[link] ({prevPtr"D"}) to[bend left] ({el"B"});%
    }
    \DoublyLinkedNode{"Z"}%
    \TailPtr{mpTail}%
    \onslide<7>{}
  \end{tikzpicture}
%\end{center}
\end{frame}

\section{Iterators}
\begin{frame}{Iterators for Container}
\begin{itemize}\setlength\itemsep{0.5em}
\item
  \emph{Iterators} are pointer-like entities that are used to\\ access individual elements in a container.
\item
  Often they are used to move sequentially from element to element, a process called \textit{iterating} through a container.
\end{itemize}
\end{frame}

\begin{frame}[fragile]{Iterators for an \texttt{\bfseries Array} Container}
  \vspace{1em}
  \begin{tikzpicture}

    % Dynamic Array
    \ArrayLabel{}
    \ArrayNodeVertical{17}
    \ArrayNodeVertical{5}
    \ArrayListVertical{23, 12, 7, 2}
    % \ArrayNodeVerticalSpan{3}
    % \ArrayNodeVertical{?}
    % \ArrayGroupVerticleRight{space}{prev}{Unused}

    % Array Class Members
    \begin{scope}[node of array/.append style={minimum width=14mm}]
      \ArrayLabel[-2.75][-2.75]{}

      \ArrayNodeStack[CSUcyan][mLength]{6}
      \node[block] (longestTitle) [fit=(index)]{};
      
      \ArrayNodeStack[CSUgold][mpArray]{}
      \fill (prev) circle(3pt);
      \coordinate (ptrCircle) at (prev);
      \draw[link] (ptrCircle) to[bend left=45] (el17);

    \end{scope}
    \begin{scope}[on background layer]
      \node[draw,fill=CSUslate, very thick,fit=(el6) (el) (longestTitle)] (stackMembers) {};
    \end{scope}
    \node[anchor=south, font=\ttfamily] at ([yshift=0pt] stackMembers.north) {\lstinline|Array<int> arr;|};

    % Iterator Members
    \onslide<2->{
      \node[pointer, name=ptr2] at ([yshift=-8.5mm] 2,-1) {};%
      \node[outer block] (last) [fit=(ptr2)]{};
      \node[label of pointer, above = of last.north west, name=head, anchor=south west, align=left, fill=codeBackground] (head) {\lstinline|Array<int>::iterator|\\
      \texttt{~~~~}\lstinline|iter = arr.begin() + 1;|};%
      \coordinate (ptr) at (ptr2);%
      \coordinate (headPtr) at (ptr2);%
      \fill (ptr) circle(3pt);%
      \draw[link] (headPtr) to[bend right=8] (el5);
    }

    \onslide<3->{
      \node[below =3pt of last.south west, anchor=north west, align=left, xshift=-3pt, execute at begin node=\setlength{\baselineskip}{1.4em}] (description) {The iterator corresponds to the \\class \lstinline|Array<int>| is of the \\type \lstinline|Array<int>::iterator|.};
    }
  \end{tikzpicture}
\end{frame}

\begin{frame}[fragile]{\texttt{\bfseries begin} and \texttt{\bfseries end} Iterator Functions}
\begin{itemize}
\item
  \lstinline{begin()}: returns an iterator to the beginning of the container.
\item
  \lstinline{end()}: returns an iterator to just past the last value.
\end{itemize}
\centerline{%
  \begin{tikzpicture}[node of array/.append style = {minimum height=7mm}]
    \ArrayLabel{\texttt{arr}}
    \ArrayNode[]{17}
    \ArrayElLabel{\lstinline|begin()|}
    \ArrayList{5, 23, 12, 7, 2}
    %\ArrayNodeSpan{5}
    %\ArrayList{,,,,,,}
    \begin{scope}[node of array/.append style={densely dotted}]
      \ArrayNode[none]{}
      \node[anchor=east, inner sep=1pt] (txtLabel) at ([shift={(0, 1)}]val.north) {Past the last element};
      \draw[->, very thick] (txtLabel.east) to[bend left=50] (val);
    \end{scope}
    \ArrayElLabel{\lstinline|end()|}
    %\node Past-the-last element
  \end{tikzpicture}
}
\end{frame}

\begin{frame}{Benefits of Iterators}
\begin{itemize}
\item<2->
  \emph{Convenience in Programming}:
  \begin{itemize}
  \item
    Provide common usage for any container (the internal\\ structure of a container does not matter).
  \item
    Provide read and/or write access to each element's values.
  \end{itemize}
\item<3->
  \emph{Reusable Code}:\\
  Iterator algorithms are not dependent on the container type.
\item<4->
  \emph{Dynamic Processing}:\\
  Easily dynamically allocate as well as delete the memory.
\end{itemize}
\end{frame}

\section*{Homework}

\begin{frame}{Homework}\Large
\begin{itemize}
\item Lab 9: \textit{Doubly-Linked Lists}
\item Project 1: \textit{Large Map}
\end{itemize}
\end{frame}

\end{document}
